
\subsubsection{2.1 Uncertainty Quantification in Neural Networks}

Calibration research established that modern neural networks are often overconfident \citep{guo2017calibration}. For language models, calibration became critical as models grew capable of generating plausible but factually incorrect text. Malinin and Gales [2018] introduced predictive uncertainty decomposition, distinguishing epistemic from aleatoric uncertainty. These methods assumed classification settings and did not address free-form generation or factuality.

\subsubsection{2.2 Token-Level Uncertainty for LLMs}

Token-level methods extract uncertainty from output logits: entropy over the vocabulary distribution, maximum probability of the selected token, or perplexity. Kadavath et al. [2022] introduced P(True), using the probability of ''True'' when asking models to self-assess as a confidence signal. These methods share a computational advantage: single-pass inference without additional sampling. However, they operate at the token level and may miss semantic-level consistency issues.

\subsubsection{2.3 Semantic-Level Uncertainty Methods}

Kuhn et al. [2023] proposed semantic entropy to capture meaning-level uncertainty. The method generates multiple samples, clusters semantically equivalent responses using NLI-based entailment, and computes entropy over cluster assignments. On custom QA evaluation, semantic entropy outperformed token entropy. Manakul et al. [2023] introduced SelfCheckGPT, using self-consistency across samples without reference documents. Both methods require 5--20x more inference compute than token-level methods.

\subsubsection{2.4 Hallucination Detection Benchmarks}

TruthfulQA \citep{lin2022truthfulqa} provides questions designed to elicit imitative falsehoods---incorrect answers humans might believe. HaluEval \citep{li2023halueval} extends to QA, dialogue, and summarization domains. These benchmarks provide ground-truth labels for hallucination detection evaluation.

\subsubsection{2.5 Methodological Gap}

Each prior method was evaluated under favorable conditions on different benchmarks. No study has compared methods on identical conditions or examined whether effectiveness depends on task format.

